\subsection{Operator Classification}

\begin{perspective} \label{per:operator-module-similarity-dictionary} \hlt{Operator Modules and Similarity}\\
Let $0<\dim\mathcal{V}<\infty$ and $S,T\in\LL(\mathcal{V})$. Then $S\sim T\Longleftrightarrow\mathcal{V}_S\cong_{\mathbb{F}[z]}\mathcal{V}_T$ by Theorem \ref{thm:operator-similarity-iff-module-isomorphism}.\\
A cyclic summand of order $f$ is isomorphic to $\mathbb{F}[z]/\langle f\rangle$, and its operator restriction has matrix $C_f$ in a cyclic basis, by Theorem \ref{thm:cyclic-subspace-quotient-and-companion}.
\end{perspective}

\begin{theorem} \label{thm:elementary-divisors-complete-similarity-invariant} \hlt{Complete Similarity Invariants}\\
Let $0<\dim\mathcal{V}<\infty$, let $S,T\in\LL(\mathcal{V})$, let $n\geq1$, and let $A,B\in\operatorname{Mat}(n,n,\mathbb{F})$.
\begin{enumerate}[label=(\roman*)]
\setlength{\itemsep}{0pt}
\item $T\sim S\Longleftrightarrow\mathcal{V}_T\cong_{\mathbb{F}[z]}\mathcal{V}_S\Longleftrightarrow\operatorname{ElemDiv}(T)=\operatorname{ElemDiv}(S)\Longleftrightarrow\operatorname{InvFact}(T)=\operatorname{InvFact}(S)$.
\item $A\sim_{\mathbb{F}}B\Longleftrightarrow(\mathbb{F}^n)_A\cong_{\mathbb{F}[z]}(\mathbb{F}^n)_B\Longleftrightarrow\operatorname{ElemDiv}(A)=\operatorname{ElemDiv}(B)\Longleftrightarrow\operatorname{InvFact}(A)=\operatorname{InvFact}(B)$.
\item If $\mathcal{B}$ is a basis of $\mathcal{V}$ and $A=[T]_{\mathcal{B}\leftarrow\mathcal{B}}$, then $\operatorname{ElemDiv}(A)=\operatorname{ElemDiv}(T)$ and $\operatorname{InvFact}(A)=\operatorname{InvFact}(T)$.
\end{enumerate}
\begin{proof}
\begin{enumerate}[label=(\roman*)]
\setlength{\itemsep}{0pt}
\item[$(i)$] ($\Rightarrow$) Suppose $T\sim S$. Then $\mathcal{V}_T\cong_{\mathbb{F}[z]}\mathcal{V}_S$ by Theorem \ref{thm:operator-similarity-iff-module-isomorphism}.\\
Both divisor data sets agree by Theorems \ref{thm:primary-cyclic-decomposition-pid}(ii) and \ref{thm:pid-invariant-factor-decomposition}(ii) with Definition \ref{def:operator-invariant-factors-elementary-divisors}(ii)--(iii).
\item[$(i)$] ($\Leftarrow$) Suppose either $\operatorname{ElemDiv}(T)=\operatorname{ElemDiv}(S)$ or $\operatorname{InvFact}(T)=\operatorname{InvFact}(S)$. Then $\mathcal{V}_T\cong_{\mathbb{F}[z]}\mathcal{V}_S$ by Theorem \ref{thm:primary-cyclic-decomposition-pid}(ii) or Theorem \ref{thm:pid-invariant-factor-decomposition}(ii), respectively, with Definition \ref{def:operator-invariant-factors-elementary-divisors}(ii)--(iii).\\
Hence $T\sim S$ and $\operatorname{ElemDiv}(T)=\operatorname{ElemDiv}(S)$ by Theorems \ref{thm:operator-similarity-iff-module-isomorphism} and \ref{thm:primary-cyclic-decomposition-pid}(ii) with Definition \ref{def:operator-invariant-factors-elementary-divisors}(iii).
\item[$(ii)$] ($\Rightarrow$) Suppose $A\sim_{\mathbb{F}}B$. Then $(\mathbb{F}^n)_A\cong_{\mathbb{F}[z]}(\mathbb{F}^n)_B$ by Corollary \ref{cor:matrix-similarity-iff-matrix-modules-isomorphic}.\\
Both divisor data sets agree by Theorems \ref{thm:primary-cyclic-decomposition-pid}(ii) and \ref{thm:pid-invariant-factor-decomposition}(ii) with Definition \ref{def:operator-invariant-factors-elementary-divisors}(ii)--(iii).
\item[$(ii)$] ($\Leftarrow$) Suppose either $\operatorname{ElemDiv}(A)=\operatorname{ElemDiv}(B)$ or $\operatorname{InvFact}(A)=\operatorname{InvFact}(B)$.\\
Then $(\mathbb{F}^n)_A\cong_{\mathbb{F}[z]}(\mathbb{F}^n)_B$ by Theorem \ref{thm:primary-cyclic-decomposition-pid}(ii) or Theorem \ref{thm:pid-invariant-factor-decomposition}(ii), respectively, with Definition \ref{def:operator-invariant-factors-elementary-divisors}(ii)--(iii).\\
Hence $A\sim_{\mathbb{F}}B$ and $\operatorname{ElemDiv}(A)=\operatorname{ElemDiv}(B)$ by Corollary \ref{cor:matrix-similarity-iff-matrix-modules-isomorphic}, Theorem \ref{thm:primary-cyclic-decomposition-pid}(ii), and Definition \ref{def:operator-invariant-factors-elementary-divisors}(iii).
\item[$(iii)$] Set $\Phi(\mathbf{v})\defeq[\mathbf{v}]_{\mathcal{B}}$. The map $\Phi$ is bijective, and $(\forall\,(p,\mathbf{v})\in\mathbb{F}[z]\times\mathcal{V})\,\Phi(p(T)\mathbf{v})=p(A)\Phi(\mathbf{v})$ by $A=[T]_{\mathcal{B}\leftarrow\mathcal{B}}$.\\
Thus $\Phi:\mathcal{V}_T\to(\mathbb{F}^n)_A$ is an $\mathbb{F}[z]$-isomorphism by Definitions \ref{def:operator-module-and-matrix-module} and \ref{def:module-homomorphism}.\\
Both divisor data sets agree by Theorems \ref{thm:primary-cyclic-decomposition-pid}(ii) and \ref{thm:pid-invariant-factor-decomposition}(ii) with Definition \ref{def:operator-invariant-factors-elementary-divisors}(ii)--(iii).
\end{enumerate}
\hlt{Structural Commentary}: Similarity identifies the operator module; PID uniqueness fixes both divisor data sets.\\
\hlt{Key Insight}: Either divisor datum determines the module and its similarity class.
\end{proof}
\end{theorem}

\begin{definition} \label{def:rational-canonical-block-forms} \hlt{Rational Canonical Block Forms}\\
Let $\mathcal{P}$ be the set of monic nonconstant polynomials in $\mathbb{F}[z]$.\\
Let $\mathcal{Q}$ be the set of positive powers of monic irreducible polynomials in $\mathbb{F}[z]$. Let $n\geq1$ and $A\in\operatorname{Mat}(n,n,\mathbb{F})$.
\begin{enumerate}[label=(\roman*)]
\setlength{\itemsep}{0pt}
\item Matrix $A$ is in \hlt{invariant-factor rational canonical form} if $(\exists\,r\in\N)(\exists\,(s_1,\ldots,s_r)\in\mathcal{P}^r)\,[A=\operatorname{diag}(C_{s_1},\ldots,C_{s_r})\land s_r\mid\cdots\mid s_1]$.
\item Matrix $A$ is in \hlt{elementary-divisor rational canonical form} if $(\exists\,m\in\N)(\exists\,(q_1,\ldots,q_m)\in\mathcal{Q}^m)\,A=\operatorname{diag}(C_{q_1},\ldots,C_{q_m})$.
\end{enumerate}
\end{definition}

\begin{theorem} \label{thm:rational-canonical-form-skeleton} \hlt{Rational Canonical Form Skeleton}\\
Let $n\geq1$, let $A\in\operatorname{Mat}(n,n,\mathbb{F})$, and write $\operatorname{InvFact}(A)=(s_1,\ldots,s_r)$. Set $R_A\defeq\operatorname{diag}(C_{s_1},\ldots,C_{s_r})$. Then:
\begin{enumerate}[label=(\roman*)]
\setlength{\itemsep}{0pt}
\item $r\in\N$, the $s_j$ are monic and nonconstant, and $s_r\mid\cdots\mid s_1$.
\item $A\sim_{\mathbb{F}}R_A$.
\item $\mu_A=s_1$ and $p_A=\prod\limits_{j=1}^rs_j$.
\end{enumerate}
\begin{proof}
\begin{enumerate}[label=(\roman*)]
\setlength{\itemsep}{0pt}
\item The tuple properties hold by Definition \ref{def:operator-invariant-factors-elementary-divisors}(ii).
\item Let $T_A:\mathbb{F}^n\to\mathbb{F}^n$, $\mathbf{x}\mapsto A\mathbf{x}$. There is a basis $\mathcal{B}$ with $[T_A]_{\mathcal{B}\leftarrow\mathcal{B}}=R_A$ by Theorem \ref{thm:operator-invariant-factor-skeleton}(ii). Hence $A\sim_{\mathbb{F}}R_A$ by Corollary \ref{cor:similar-iff-same-linear-operator}.
\item Identities $\mu_{T_A}=s_1$ and $q_{T_A}=\prod\limits_{j=1}^rs_j$ hold by Theorem \ref{thm:operator-invariant-factor-skeleton}(iii)--(iv). Thus $\mu_A=s_1$, $p_A=\prod\limits_{j=1}^rs_j$ by Definitions \ref{def:minimal-polynomial-matrix}, \ref{def:operator-characteristic-polynomial}.
\end{enumerate}
\hlt{Structural Commentary}: The invariant-factor decomposition supplies a cyclic basis whose coordinate blocks represent $A$.\\
\hlt{Key Insight}: The divisibility-ordered tuple fixes the companion-block representative.
\end{proof}
\end{theorem}

\begin{theorem} \label{thm:rational-canonical-form-elementary-and-invariant-versions} \hlt{Rational Canonical Form: Two Versions}\\
Let $0<\dim\mathcal{V}<\infty$, let $T\in\LL(\mathcal{V})$, let $n\geq1$, and let $A\in\operatorname{Mat}(n,n,\mathbb{F})$.
\begin{enumerate}[label=(\roman*)]
\setlength{\itemsep}{0pt}
\item If $\operatorname{ElemDiv}(T)=\{q_1,\ldots,q_m\}$, then $(\exists\,\mathcal{R})\,[T]_{\mathcal{R}\leftarrow\mathcal{R}}=\operatorname{diag}(C_{q_1},\ldots,C_{q_m})$ for a basis $\mathcal{R}$ of $\mathcal{V}$.
\item The elementary-divisor form in (i) is unique up to block permutation.
\item If $\operatorname{InvFact}(T)=(s_1,\ldots,s_r)$, then $(\exists\,\mathcal{B})\,[T]_{\mathcal{B}\leftarrow\mathcal{B}}=\operatorname{diag}(C_{s_1},\ldots,C_{s_r})$ for a basis $\mathcal{B}$ of $\mathcal{V}$.
\item The invariant-factor form in (iii) is unique under the order $s_r\mid\cdots\mid s_1$.
\item The class $[A]_{\sim_{\mathbb{F}}}$ contains both forms, with the uniqueness in (ii) and (iv).
\end{enumerate}
\begin{proof}
\begin{enumerate}[label=(\roman*)]
\setlength{\itemsep}{0pt}
\item Basis $\mathcal{R}$ with displayed block matrix is obtained by concatenating cyclic bases in Theorem \ref{thm:operator-primary-cyclic-decomposition}(ii)--(iii) and Theorem \ref{thm:cyclic-subspace-quotient-and-companion}(iii).
\item Each block $C_q$ is cyclic with order $q$ by Definition \ref{def:companion-matrix}, Theorem \ref{thm:companion-matrix-has-prescribed-minimal-polynomial}(i), and Proposition \ref{prop:minimal-polynomial-module-order}(i).\\
Hence an elementary-divisor block matrix has the multiset of its block orders by Proposition \ref{prop:elementary-divisors-direct-sums} and Definition \ref{def:operator-invariant-factors-elementary-divisors}(iii).\\
Every coordinate matrix of $T$ has divisor data equal to those of $T$ by Theorem \ref{thm:elementary-divisors-complete-similarity-invariant}(iii), so only the block order can vary.
\item A basis $\mathcal{B}$ with the displayed block matrix exists by Theorem \ref{thm:operator-invariant-factor-skeleton}(i)--(ii).
\item The companion blocks in an invariant-factor form are cyclic with their defining orders by Definition \ref{def:companion-matrix} and Theorem \ref{thm:companion-matrix-has-prescribed-minimal-polynomial}(i).\\
Their divisibility-ordered list is the invariant-factor tuple of the coordinate module by Theorem \ref{thm:pid-invariant-factor-decomposition}(i)--(ii).\\
It equals $\operatorname{InvFact}(T)$ by Theorem \ref{thm:elementary-divisors-complete-similarity-invariant}(iii), so the list is unique.
\item Both forms for $T_A:\mathbb{F}^n\to\mathbb{F}^n$, $\mathbf{x}\mapsto A\mathbf{x}$, exist with the stated uniqueness by parts (i)--(iv).\\
Each resulting coordinate matrix is similar to $A$ by Corollary \ref{cor:similar-iff-same-linear-operator}.
\end{enumerate}
\hlt{Structural Commentary}: Primary and invariant-factor decompositions concatenate cyclic bases in two intrinsic block orders.\\
\hlt{Key Insight}: The two forms preserve one module while exposing different divisor groupings.
\end{proof}
\end{theorem}

\begin{computationmethod} \label{meth:construct-operator-from-elementary-divisors} \hlt{Constructing an Operator from Elementary Divisors}\\
Given $m\in\N$ and a multiset $Q=\{q_1,\ldots,q_m\}$ of monic nonconstant prime powers in $\mathbb{F}[z]$, set $n\defeq\sum\limits_{j=1}^m\deg q_j$.\\
The output is a matrix in $\operatorname{Mat}(n,n,\mathbb{F})$ with elementary-divisor multiset $Q$ and its invariant-factor form.
\begin{enumerate}[label=(\roman*)]
\setlength{\itemsep}{0pt}
\item Set $R_Q\defeq\operatorname{diag}(C_{q_1},\ldots,C_{q_m})\in\operatorname{Mat}(n,n,\mathbb{F})$.
\item $(\forall\,j\in\{1,\ldots,m\})\,\mathcal{C}_{T_{C_{q_j}}}(\mathbf{e}_1)=\mathbb{F}^{\deg q_j}$ by Definition \ref{def:companion-matrix}. Identities $(\forall\,j\in\{1,\ldots,m\})\,\mu_{C_{q_j}}=q_j$ hold by Theorem \ref{thm:companion-matrix-has-prescribed-minimal-polynomial}(i).\\
These polynomials represent the corresponding module orders by Proposition \ref{prop:minimal-polynomial-module-order}(i).\\
Hence $(\forall\,j\in\{1,\ldots,m\})\,\operatorname{ElemDiv}(T_{C_{q_j}})=\{q_j\}$ by Theorem \ref{thm:cyclic-torsion-module-characterization}(i),(iii) and Definition \ref{def:operator-invariant-factors-elementary-divisors}(iii).\\
Therefore $\operatorname{ElemDiv}(T_{R_Q})=Q$ by repeated application of Proposition \ref{prop:elementary-divisors-direct-sums}.
\item Recombine $Q$ into $\operatorname{InvFact}(T_{R_Q})=(s_1,\ldots,s_r)$ by Method \ref{meth:elementary-divisors-to-invariant-factors}. Then $\operatorname{diag}(C_{s_1},\ldots,C_{s_r})\sim_{\mathbb{F}}R_Q$ by Theorem \ref{thm:rational-canonical-form-skeleton}(ii).
\end{enumerate}
\hlt{Key Insight}: Companion blocks realise cyclic summands; block sums realise direct sums.
\end{computationmethod}

\begin{example} \label{ex:rational-canonical-forms-from-minimal-polynomial} \hlt{Rational Canonical Possibilities from Degree Counts}\\
Let $T\in\LL(\mathbb{R}^7)$ satisfy $\mu_T=(z-1)(z^2+1)^2$, and set $Q\defeq\operatorname{ElemDiv}(T)$.\\
Then $\operatorname{lcm}Q=\mu_T$ and $\sum\limits_{q\in Q}\deg q=7$ by Theorem \ref{thm:minimal-characteristic-elementary-divisor-relations}(i),(iii) and Definition \ref{def:characteristic-polynomial-from-elementary-divisors}.\\
Let $a,b,c$ be the multiplicities in $Q$ of $z-1$, $z^2+1$, and $(z^2+1)^2$, respectively.\\
The constraints are $a,c\geq1$, $b\geq0$, and $a+2b+4c=7$, so $c=1$ and $(a,b)\in\{(1,1),(3,0)\}$.\\
Thus $Q_1=\{z-1,z^2+1,(z^2+1)^2\}$ or $Q_2=\{z-1,z-1,z-1,(z^2+1)^2\}$.\\
The two possible elementary-divisor forms are $R_i=\bigoplus\limits_{q\in Q_i}C_q$ $(i\in\{1,2\})$, up to block order.
\end{example}

\begin{theorem} \label{thm:determinant-form-characteristic-polynomial} \hlt{Determinant Form of the Characteristic Polynomial}\\
Let $n\geq1$ and $A\in\operatorname{Mat}(n,n,\mathbb{F})$.
\begin{enumerate}[label=(\roman*)]
\setlength{\itemsep}{0pt}
\item $\chi_A=p_A$, so $\chi_A(z)=\det(zI_n-A)$.
\item If $0<\dim\mathcal{V}=n$, $T\in\LL(\mathcal{V})$, and $A=[T]_{\mathcal{B}\leftarrow\mathcal{B}}$ for a basis $\mathcal{B}$, then $\chi_T=q_T=\chi_A$.
\end{enumerate}
\begin{proof}
\begin{enumerate}[label=(\roman*)]
\setlength{\itemsep}{0pt}
\item $\chi_A=\chi_{T_A}$ by Definition \ref{def:characteristic-polynomial-from-elementary-divisors}. $\chi_{T_A}=q_{T_A}=p_A$ and $p_A(z)=\det(zI_n-A)$ hold by Theorem \ref{thm:minimal-characteristic-elementary-divisor-relations}(ii), Definitions \ref{def:operator-characteristic-polynomial}, \ref{def:matrix-characteristic-polynomial}.
\item The identity $\chi_T=q_T=p_A$ holds by Theorem \ref{thm:minimal-characteristic-elementary-divisor-relations}(ii) and Definition \ref{def:operator-characteristic-polynomial}. Hence $p_A=\chi_A$ by part (i).
\end{enumerate}
\hlt{Structural Commentary}: Elementary-divisor products equal operator polynomial, which coordinates identify with shifted determinant.\\
\hlt{Key Insight}: The abstract block product is the concrete characteristic determinant.
\end{proof}
\end{theorem}

\begin{theorem} \label{thm:base-field-change-elementary-divisors} \hlt{Changing the Base Field}\\
Let $\mathbb{F}\subseteq\mathbb{K}$ be fields, $n\geq1$, $A\in\operatorname{Mat}(n,n,\mathbb{F})$. Subscripts on $\operatorname{ElemDiv}$ and $\operatorname{InvFact}$ indicate the coefficient field.\\
Let $s\in\N$, $(k_i)_{i=1}^s\in\N^s$, $p_i\in\mathbb{F}[z]$ be distinct monic irreducibles; $\operatorname{ElemDiv}_{\mathbb{F}}(A)=\{p_i^{e_{i,j}}:1\leq i\leq s,\ 1\leq j\leq k_i\}$, where $e_{i,j}\in\N$.\\
For $i\in\{1,\ldots,s\}$, let $m_i\in\N$, $(d_{i,\ell})_{\ell=1}^{m_i}\in\N^{m_i}$, $a_{i,\ell}\in\mathbb{K}[z]$ be pairwise distinct monic irreducibles. Suppose $p_i=\prod\limits_{\ell=1}^{m_i}a_{i,\ell}^{d_{i,\ell}}$. Then:
\begin{enumerate}[label=(\roman*)]
\setlength{\itemsep}{0pt}
\item $\operatorname{ElemDiv}_{\mathbb{K}}(A)=\{a_{i,\ell}^{d_{i,\ell}e_{i,j}}:1\leq i\leq s,\ 1\leq j\leq k_i,\ 1\leq\ell\leq m_i\}$ as multisets;
\item $\operatorname{InvFact}_{\mathbb{K}}(A)=\operatorname{InvFact}_{\mathbb{F}}(A)$ under the coefficient inclusion $\mathbb{F}[z]\subseteq\mathbb{K}[z]$.
\end{enumerate}
\begin{proof}
\begin{enumerate}[label=(\roman*)]
\setlength{\itemsep}{0pt}
\item Set $R_{\mathbb{F}}\defeq\bigoplus\limits_{i=1}^s\bigoplus\limits_{j=1}^{k_i}C_{p_i^{e_{i,j}}}$. The similarities $A\sim_{\mathbb{F}}R_{\mathbb{F}}$ and $A\sim_{\mathbb{K}}R_{\mathbb{F}}$ hold by Theorem \ref{thm:rational-canonical-form-elementary-and-invariant-versions}(v) and $\operatorname{GL}_n(\mathbb{F})\subseteq\operatorname{GL}_n(\mathbb{K})$.\\
Fix $(i,j)$. The first standard vector is cyclic for $C_{p_i^{e_{i,j}}}$ by Definition \ref{def:companion-matrix}.\\
Its order over $\mathbb{K}[z]$ is $p_i^{e_{i,j}}$ by Theorem \ref{thm:companion-matrix-has-prescribed-minimal-polynomial}(i) and Proposition \ref{prop:minimal-polynomial-module-order}(i).\\
The factors $a_{i,\ell}^{d_{i,\ell}e_{i,j}}$ are pairwise coprime, so $C_{p_i^{e_{i,j}}}\sim_{\mathbb{K}}\bigoplus\limits_{\ell=1}^{m_i}C_{a_{i,\ell}^{d_{i,\ell}e_{i,j}}}$ by Theorems \ref{thm:cyclic-composition-coprime-orders}(iii) and \ref{thm:cyclic-subspace-quotient-and-companion}(iii).\\
Replacing every block defines an elementary-divisor form $R_{\mathbb{K}}$ with $A\sim_{\mathbb{K}}R_{\mathbb{K}}$ by the direct sum of the blockwise similarity matrices.\\
Each block of $R_{\mathbb{K}}$ is cyclic with its displayed order by Definition \ref{def:companion-matrix}, Theorem \ref{thm:companion-matrix-has-prescribed-minimal-polynomial}(i), and Proposition \ref{prop:minimal-polynomial-module-order}(i).\\
Thus the elementary divisors of $R_{\mathbb{K}}$ are its block orders by Proposition \ref{prop:elementary-divisors-direct-sums} and Definition \ref{def:operator-invariant-factors-elementary-divisors}(iii).\\
The asserted multiset is therefore $\operatorname{ElemDiv}_{\mathbb{K}}(A)$ by Theorem \ref{thm:elementary-divisors-complete-similarity-invariant}(ii).
\item Write $\operatorname{InvFact}_{\mathbb{F}}(A)=(s_1,\ldots,s_r)$. Then $A\sim_{\mathbb{F}}\operatorname{diag}(C_{s_1},\ldots,C_{s_r})$ by Theorem \ref{thm:rational-canonical-form-skeleton}(ii).\\
The same similarity holds over $\mathbb{K}$, and the $s_j$ remain monic and nonconstant with $s_r\mid\cdots\mid s_1$, by coefficient inclusion.\\
This block matrix is invariant-factor form over $\mathbb{K}$ by Definition \ref{def:rational-canonical-block-forms}(i); claimed tuple equality holds by Theorem \ref{thm:rational-canonical-form-elementary-and-invariant-versions}(iv)--(v).
\end{enumerate}
\hlt{Structural Commentary}: Coprime factorization refines each elementary block; the original divisibility-ordered form remains valid.\\
\hlt{Key Insight}: Scalar extension splits prime powers without changing invariant-factor blocks.
\end{proof}
\end{theorem}

\begin{corollary} \label{cor:nonderogatory-iff-companion-similar} \hlt{Nonderogatory Means One Companion Block}\\
Let $n\geq1$ and $A\in\operatorname{Mat}(n,n,\mathbb{F})$. The following are equivalent:
\begin{enumerate}[label=(\roman*)]
\setlength{\itemsep}{0pt}
\item $A$ is nonderogatory;
\item $(\mathbb{F}^n)_A$ is cyclic;
\item $A\sim_{\mathbb{F}}C_{\chi_A}$.
\end{enumerate}
\begin{proof}
\begin{enumerate}[label=(\roman*)]
\setlength{\itemsep}{0pt}
\item[$(i)$] $\Rightarrow (ii)$ The operator $T_A$ is nonderogatory by condition (i) and Definition \ref{def:nonderogatory-operator-module}; $(\mathbb{F}^n)_A$ is cyclic by Theorem \ref{thm:nonderogatory-cyclic-equivalences}(iv),(i).
\item[$(ii)$] $\Rightarrow (iii)$ $\exists$ $s\in\mathbb{F}[z]$ with $\operatorname{InvFact}(A)=(s)$ by condition (ii), Theorem \ref{thm:cyclic-torsion-module-characterization}(i),(iv), Theorem \ref{thm:pid-invariant-factor-decomposition}(i)--(ii), and Definition \ref{def:operator-invariant-factors-elementary-divisors}(ii).\\
The identities $A\sim_{\mathbb{F}}C_s$ and $s=q_{T_A}=p_A=\chi_A$ hold by Theorem \ref{thm:operator-invariant-factor-skeleton}(ii)--(iv), Theorem \ref{thm:determinant-form-characteristic-polynomial}(i), and Definition \ref{def:operator-characteristic-polynomial}.
\item[$(iii)$] $\Rightarrow (i)$ The identity $\mu_A=\mu_{C_{\chi_A}}=\chi_A$ holds by condition (iii), Theorem \ref{thm:companion-matrix-has-prescribed-minimal-polynomial}(i), and Proposition \ref{prop:minimal-polynomial-module-order}(iii).\\
Thus $A$ is nonderogatory by Theorem \ref{thm:determinant-form-characteristic-polynomial}(i) and Definitions \ref{def:operator-characteristic-polynomial} and \ref{def:nonderogatory-operator-module}.
\end{enumerate}
\hlt{Structural Commentary}: Cyclic-module characterization reduces nonderogatoriness to one invariant-factor companion block.\\
\hlt{Key Insight}: A cyclic generator produces a single companion block.
\end{proof}
\end{corollary}

\begin{corollary} \label{cor:diagonal-rational-canonical-form-criterion} \hlt{When the Elementary-Divisor Form Is Diagonal}\\
Let $0<\dim\mathcal{V}<\infty$, let $T\in\LL(\mathcal{V})$, set $Q\defeq\operatorname{ElemDiv}(T)$, and set $R_{\operatorname{ed}}\defeq\bigoplus\limits_{q\in Q}C_q$. The following are equivalent:
\begin{enumerate}[label=(\roman*)]
\setlength{\itemsep}{0pt}
\item $R_{\operatorname{ed}}$ is diagonal;
\item $(\forall\,q\in Q)\,\deg q=1$;
\item $(\exists\,\Lambda\subseteq\mathbb{F})\,[0<\abs{\Lambda}<\infty\land\mu_T=\prod\limits_{\lambda\in\Lambda}(z-\lambda)]$.
\end{enumerate}
\begin{proof}
\begin{enumerate}[label=(\roman*)]
\setlength{\itemsep}{0pt}
\item[$(i)$] $\Rightarrow (ii)$ Every companion block $C_q$ is diagonal by condition (i).\\
Hence $(\forall\,q\in Q)\,\deg q=1$ by the nonzero subdiagonal in Definition \ref{def:companion-matrix} whenever $\deg q\geq2$.
\item[$(ii)$] $\Rightarrow (iii)$ Polynomial $\mu_T$ is least common multiple of monic linear polynomials in $Q$ by condition (ii) and Theorem \ref{thm:minimal-characteristic-elementary-divisor-relations}(i).\\
Hence condition (iii) holds for the finite nonempty set of their roots.
\item[$(iii)$] $\Rightarrow (i)$ The divisibility $(\forall\,q\in Q)\,q\mid\mu_T$ holds by condition (iii) and Theorem \ref{thm:minimal-characteristic-elementary-divisor-relations}(i).\\
Since the elements of $Q$ are prime powers and $\mu_T$ is square-free, $(\forall\,q\in Q)(\exists\,\lambda\in\Lambda)\,q=z-\lambda$.\\
Thus every $C_q$ is a $1\times1$ diagonal matrix by Definition \ref{def:companion-matrix}, so $R_{\operatorname{ed}}$ is diagonal.
\end{enumerate}
\hlt{Structural Commentary}: Companion matrices of degree at least two have a nonzero subdiagonal; $\mu_T$ detects the remainder.\\
\hlt{Key Insight}: Diagonal companion blocks are exactly the linear elementary divisors.
\end{proof}
\end{corollary}

\begin{example} \label{ex:base-field-rational-canonical-form-minimal-polynomial} \hlt{The Same Polynomial over Different Fields}\\
Set $\mu\defeq(z^2+1)(z^2-2)$, $A\defeq C_{\mu}\in\operatorname{Mat}(4,4,\mathbb{Q})$. Then $\mu_A=p_A=\mu$ by Theorem \ref{thm:companion-matrix-has-prescribed-minimal-polynomial}(i)--(ii).\\
Hence $A$ is nonderogatory by Definitions \ref{def:operator-characteristic-polynomial} and \ref{def:nonderogatory-operator-module}.\\
The matrix $A=C_\mu$ is an invariant-factor form by Definition \ref{def:rational-canonical-block-forms}(i), so $\operatorname{InvFact}_{\mathbb{Q}}(A)=(\mu)$ by Theorem \ref{thm:rational-canonical-form-elementary-and-invariant-versions}(iv)--(v).\\
Thus the invariant-factor form over $K$ is $C_\mu$ for $K\in\{\mathbb{Q},\mathbb{R},\mathbb{C}\}$ by Theorem \ref{thm:base-field-change-elementary-divisors}(ii). The elementary divisors are:
\begin{center}
\begin{tabular}{c|c}
$K$ & $\operatorname{ElemDiv}_K(A)$\\
\hline
$\mathbb{Q}$ & $\{z^2+1,\ z^2-2\}$\\
$\mathbb{R}$ & $\{z^2+1,\ z-\sqrt{2},\ z+\sqrt{2}\}$\\
$\mathbb{C}$ & $\{z-i,\ z+i,\ z-\sqrt{2},\ z+\sqrt{2}\}$
\end{tabular}
\end{center}
The elementary-divisor form over $K$ is $\bigoplus\limits_{q\in\operatorname{ElemDiv}_K(A)}C_q$.\\
Scalar extension refines these blocks while preserving $C_\mu$; over $\mathbb{C}$, the elementary-divisor form is diagonal.
\end{example}

\begin{perspective} \label{per:module-and-spectral-canonical-form-crosswalk} \hlt{Crosswalk with the Spectral Presentation}\\
Let $0<\dim_{\mathbb{C}}\mathcal{V}<\infty$ and $T\in\LL(\mathcal{V})$. The module and spectral presentations encode the same data at different splitting levels:
\begin{enumerate}[label=(\roman*)]
\setlength{\itemsep}{0pt}
\item $T$ is diagonalizable exactly when $(\forall\,q\in\operatorname{ElemDiv}(T))\,\deg q=1$ by Theorem \ref{thm:diagonalizable-iff-minimal-polynomial-squarefree} and Corollary \ref{cor:diagonal-rational-canonical-form-criterion};
\item the elementary divisor $(z-\lambda)^k$ corresponds to the Jordan block $J_k(\lambda)$ by Theorem \ref{thm:complex-jordan-form} and Perspective \ref{per:elementary-divisors-and-jordan-blocks};
\item the rank drops of the $(z-\lambda)$-primary component recover the block exponents $k$ by Proposition \ref{prop:p-primary-layer-counts-recover-exponents}.
\end{enumerate}
\hlt{Key Insight}: Rational form is Jordan form before irreducibles split.
\end{perspective}

\begin{computationmethod} \label{meth:operator-module-to-spectral-form} \hlt{Moving Between Module and Spectral Data}\\
Let $\mathbb{F}\subseteq\mathbb{C}$, $0<\dim_{\mathbb{F}}\mathcal{V}<\infty$, $T\in\LL(\mathcal{V})$. Choose basis $\mathcal{B}$, set $A\defeq[T]_{\mathcal{B}\leftarrow\mathcal{B}}$, and let $\mathbb{K}\subseteq\mathbb{C}$ be the splitting field of $\mu_T$ over $\mathbb{F}$.\\
The outputs are a rational canonical form over $\mathbb{F}$ and the Jordan-block data of $A$ over $\mathbb{K}$.
\begin{enumerate}[label=(\roman*)]
\setlength{\itemsep}{0pt}
\item Factor $\mu_T=\prod\limits_{i=1}^s p_i^{e_i}$ into distinct monic irreducibles over $\mathbb{F}$, and set $\mathcal{V}_{p_i}\defeq\ker p_i(T)^{e_i}$.
\item Decompose each $\mathcal{V}_{p_i}$ into cyclic summands of orders $p_i^{e_{i,j}}$ by Theorem \ref{thm:operator-primary-cyclic-decomposition}(ii).\\
These orders form $\operatorname{ElemDiv}(T)$ by part (iii) of that theorem.
\item Assemble the blocks $C_{p_i^{e_{i,j}}}$ into an elementary-divisor rational canonical form over $\mathbb{F}$ by Theorem \ref{thm:rational-canonical-form-elementary-and-invariant-versions}(i).
\item Factor the $p_i$ over $\mathbb{K}$ and refine the elementary divisors into powers $(z-\lambda)^k$ by Theorem \ref{thm:base-field-change-elementary-divisors}(i) for $A$.\\
In the reversed shifted-power basis of $\mathbb{K}[z]/\langle(z-\lambda)^k\rangle$, multiplication by $z$ has matrix $J_k(\lambda)$.\\
Hence each refined cyclic block is similar over $\mathbb{K}$ to $J_k(\lambda)$ by Theorem \ref{thm:cyclic-subspace-quotient-and-companion}(i),(iii).
\item If $T$ is diagonalizable over $\mathbb{F}$, a diagonal coordinate matrix is a direct sum of blocks $C_{z-\lambda}$.\\
Elementary divisors linear by Proposition \ref{prop:elementary-divisors-direct-sums}, Definition \ref{def:operator-invariant-factors-elementary-divisors}(iii); $(\forall\,q\in\operatorname{ElemDiv}(T))\,\deg q=1$ by Theorem \ref{thm:elementary-divisors-complete-similarity-invariant}(iii).\\
Conversely, if every elementary divisor has degree one, form in (iii) is diagonal by Corollary \ref{cor:diagonal-rational-canonical-form-criterion}; $T$ is diagonalizable over $\mathbb{F}$.
\end{enumerate}
\hlt{Key Insight}: Rational and Jordan forms encode the same elementary-divisor data over different fields.
\end{computationmethod}

Figure \ref{fig:operator-module-canonical-form-dictionary} tracks a similarity class through its operator module, divisor data, and companion-block representative.

\begin{illustrationfigure}[H]
  \centering
  \begin{tikzpicture}[
  >=Stealth,
  every node/.style={font=\small},
  class/.style={draw=blue!60!black, fill=blue!7, rounded corners=2pt,
    minimum width=2.25cm, minimum height=0.72cm, align=center, inner sep=5pt},
  data/.style={draw=dartmouthgreen!65!black, fill=dartmouthgreen!8, rounded corners=2pt,
    minimum width=2.9cm, minimum height=0.72cm, align=center, inner sep=5pt},
  form/.style={draw=gray!60, fill=gray!6, rounded corners=2pt,
    minimum width=3.45cm, minimum height=0.72cm, align=center, inner sep=5pt},
  link/.style={->, line width=0.75pt, draw=black!70},
  note/.style={font=\scriptsize, inner sep=2pt, fill=white}
]
  \node[class] (ops) at (0,0) {$[T]_{\sim}$};
  \node[class] (mods) at (4.3,0) {$[\mathcal{V}_T]_{\mathbb{F}[z]}$};
  \node[data] (invs) at (8.7,0) {$\{s_i\}:s_r\mid\cdots\mid s_1$\\$\{q_j\}$};
  \node[form] (mats) at (13.4,0) {$\operatorname{diag}(C_{s_1},\ldots,C_{s_r})$};

  \draw[link] (ops) -- node[note, above] {$p\cdot\mathbf{v}=p(T)\mathbf{v}$} (mods);
  \draw[link] (mods) -- node[note, above] {$\bigoplus_i\mathbb{F}[z]/\langle s_i\rangle$} (invs);
  \draw[link] (invs) -- node[note, above] {$s_i\mapsto C_{s_i}$} (mats);
\end{tikzpicture}

  \caption{Similarity classes pass through operator modules and divisor data to companion-block representatives.}
  \label{fig:operator-module-canonical-form-dictionary}
\end{illustrationfigure}
